\chapter*{Abstract}

This thesis presents the design and validation of a custom analog front-end for simultaneous acquisition of electromyography (EMG) and electrodermal activity (EDA) from the same muscle site. The target application is objective stress and pain monitoring in non-verbal patients, for whom palmar EDA placement is impractical and somatic-autonomic co-localization at a clinically relevant muscle site is needed.

Co-locating the two modalities introduces a ground loop between the EMG and EDA circuits. If the EMG reference electrode offers a lower-impedance return path than the intended EDA reference, part of the EDA excitation current diverts through the EMG chain and corrupts the skin conductance reading. The front-end addresses this by placing the EMG and EDA signal chains in separate ground domains, with reinforced galvanic isolation between each front-end and the data acquisition unit. Power and signal isolation together break the shared current path while letting both channels be sampled on a common data acquisition card.

The system was validated through five recording sessions on a single participant. Sessions~1 and~2 characterized each front-end individually. Sessions~3 to~5 evaluated simultaneous operation, mutual interference between the two channels, and reproducibility across days. The results show that galvanic isolation supports co-located EMG and direct current EDA acquisition at the same muscle site, and the design offers a starting point for future multimodal stress and pain monitoring in non-verbal patient populations.